\subsection{ADMISSIBLE LATTICES}

Generalizing the terminology adopted for vector spaces, an endomorphism u of a module M is said to be diagonalizable if there exists a basis of M such that the matrix of u relative to this basis is diagonal.

Lemma 6. Let M be a free \(\mathbf{Z}\)-module of finite type, u an endomorphism of M, and v the endomorphism \(u \otimes 1\) of \(M \otimes_{\mathbf{Z}} \mathbf{Q}\) . Assume that \(\binom{v}{n}(M) \subset M\) for all \(n \in \mathbf{N}\) . Then u is diagonalizable.

\begin{enumerate}
\def\labelenumi{\alph{enumi})}
\item
  For any polynomial \(\mathrm{P} \in \mathbf{Q}[\mathrm{T}]\) such that \(\mathrm{P}(\mathbf{Z}) \subset \mathbf{Z}\) , we have \(\mathrm{P}(v)(\mathrm{M}) \subset \mathrm{M}\) (no. 4, Cor. of Prop. 2), so det \(\mathrm{P}(v) \in \mathbf{Z}\) .
\item
  Denote by \(\chi_{v}(t)=t^{d}+\alpha_{1}t^{d-1}+\cdots\) the characteristic polynomial of v. Let \(k\in\mathbf{Z}, n\in\mathbf{N}\) . Applying a) to the polynomial \(\binom{T-k}{n}\) , we see that the number
\end{enumerate}

\[
\begin{array}{l} a _ {n} = \det \binom {v - k} {n} = \frac {1}{(n !) ^ {d}} \det (v - k) \det (v - k - 1) \dots \det (v - k - n + 1) \\ = \frac {(- 1) ^ {n}}{(n !) ^ {4}} \chi_ {v} (k) \chi_ {v} (k + 1) \dots \chi_ {v} (k + n - 1) \end{array}
\]

is an integer. Take \(k - 1 < -\alpha_{1} / d\) . Then

\[
\chi_ {v} (k + n - 1) = n ^ {d} + (\alpha_ {1} + (k - 1) d) n ^ {d - 1} + \dots
\]

and

\[
| a _ {n} | = \frac {| \chi_ {v} (k + n - 1) |}{n ^ {d}} | a _ {n - 1} |;
\]

hence, if \(a_{n} \neq 0\) for all \(n \in \mathbf{N}\) , the sequence of the \(|a_{n}|\) is strictly decreasing for n sufficiently large, which is absurd. It follows that v has an integer eigenvalue \(\lambda\) . Put \(\mathrm{M}' = \mathrm{Ker}(u - \lambda.1)\) and \(M'' = M/M'\) . Then \(M'\) is the intersection with M of a vector subspace of \(M \otimes_{\mathbf{Z}} \mathbf{Q}\) , so the \(\mathbf{Z}\)-module \(M''\) is torsion-free of finite type, and consequently free of rank \textless{} d.~Arguing by induction on d and applying the induction hypothesis to the endomorphism of \(M''\) induced by u, we conclude that all the eigenvalues of v in an algebraically closed extension of \(\mathbf{Q}\) are integers.

\begin{enumerate}
\def\labelenumi{\alph{enumi})}
\setcounter{enumi}{2}
\tightlist
\item
  We show that \(v\) is diagonalizable. Let \(\lambda\) be an eigenvalue of \(v\) and let \(x \in \mathrm{M} \otimes_{\mathbf{Z}} \mathbf{Q}\) be such that \((v - \lambda)^2 x = 0\) . We have \(v(vx - \lambda x) = \lambda(vx - \lambda x)\) , so
\end{enumerate}

\[
\begin{array}{c} \frac {1}{n !} (v - \lambda - n + 1) (v - \lambda - n + 2) \ldots (v - \lambda - 1) (v - \lambda) x \\ = \frac {(- 1) ^ {n - 1}}{n} (v x - \lambda x). \end{array}
\]

By \(a\) ), this implies that \(vx - \lambda x \in n\mathbf{M}\) for all \(n \in \mathbf{N}\) , so \((v - \lambda)x = 0\) .

\begin{enumerate}
\def\labelenumi{\alph{enumi})}
\setcounter{enumi}{3}
\tightlist
\item
  Let \(\lambda\) be an eigenvalue of v and let \((\lambda - a, \lambda + b]\) be an interval in \(\mathbf{Z}\) containing all the eigenvalues of v. Consider the polynomial
\end{enumerate}

\[
\begin{array}{c} \mathrm{P(T)} = (- 1) ^ {b} \frac {(\mathrm{T} - \lambda - 1) (\mathrm{T} - \lambda - 2) \ldots (\mathrm{T} - \lambda - b)}{b !} \\ \times \frac {(\mathrm{T} - \lambda + 1) (\mathrm{T} - \lambda + 2) \ldots (\mathrm{T} - \lambda + a)}{a !}. \end{array}
\]

We have \(\mathrm{P}(\mathbf{Z})\subset\mathbf{Z},\mathrm{P}(\lambda)=1,\mathrm{P}(\mu)=0\) for \(\mu\in\mathbf{Z}\cap(\lambda-a,\lambda+b)\) and \(\mu\neq\lambda\) . By a), \(\mathrm{P}(v)(\mathrm{M})\subset\mathrm{M}\) . By c), \(\mathrm{P}(v)\) is a projection of \(M\otimes_{\mathbf{Z}}\mathbf{Q}\) onto the eigenspace corresponding to \(\lambda\) .

Q.E.D.

Remark 1. If we only assume that v is diagonalizable with integer eigenvalues, u is not necessarily diagonalizable (for example, take \(M = \mathbf{Z}^{2}\) and \(u(x, y) = (y, x)\) for all \((x, y) \in M\) ).

Let \(\mathfrak{g},\mathfrak{h},\mathrm{R},\mathcal{H},\mathcal{G}^{\alpha},\mathcal{G},\mathcal{U}\) be as in no. 7, and assume that \(\mathcal{G}\) is a Chevalley order in \((\mathfrak{g},\mathfrak{h})\) .

DEFINITION 3. Let E be a \(\mathfrak{g}\)-module. A lattice \(\mathcal{E}\) in E is said to be admissible (relative to \(\mathcal{G}\)) if the following conditions are satisfied:

\begin{enumerate}
\def\labelenumi{(\roman{enumi})}
\item
  \(\mathcal{U}\) maps \(\mathcal{E}\) to \(\mathcal{E}\) ;
\item
  \(\mathcal{E}\) is stable under \(\binom{h}{n}\) and \(x^{(n)}\) for all \(\alpha \in \mathrm{R}, x \in \mathcal{G}^{\alpha}, n \in \mathbf{N}, h \in \mathcal{H}\) .
\end{enumerate}

Remarks. 2) Let \(\rho\) be the adjoint representation of \(\mathfrak{g}\) on \(\mathrm{U}(\mathfrak{g})\) . Let \(\alpha, x, n, h\) be as in (ii) above. We have \(\rho(x^{(n)})\) . \(\mathcal{U} \subset \mathcal{U}\) by Lemma 2. On the other hand, if \(p \in \mathbf{N}\) ,

\[
\rho \left(\binom{h}{p}\right) x ^ {(n)} = \binom{\operatorname{ad} h}{p} x ^ {(n)} = \binom{n \alpha (h)}{p} x ^ {(n)}
\]

(no. 5, formula (13)), so \(\rho\left(\binom{h}{p}\right). \mathcal{U} \subset \mathcal{U}\) . This proves that \(\mathcal{U}\) is an admissible lattice in \(\mathrm{U}(\mathfrak{g})\) , and it follows that \(\mathcal{G}\) is an admissible lattice in \(\mathfrak{g}\) (for the adjoint representation).

\begin{enumerate}
\def\labelenumi{\arabic{enumi})}
\setcounter{enumi}{2}
\tightlist
\item
  Let E be a finite dimensional \(\mathfrak{g}\)-module, \(\mathcal{E}\) an admissible lattice in E, \(\mathfrak{c}\) the centre of \(\mathfrak{g}\). By Lemma 6, every element of \(\mathfrak{c}\) defines a diagonalizable endomorphism of E. Hence E is semi-simple (Chap. I, §6, no. 5, Th. 4). Thus, E is a direct sum of simple \(\mathcal{D}\mathfrak{g}\)-modules on which \(\mathfrak{c}\) induces homotheties. By Lemma 6, \(\mathcal{E}=\oplus(\mathcal{E}\cap\mathrm{E}^{\lambda})\) and, for all weights \(\lambda\) of E, we have
\end{enumerate}

\[
\lambda (\mathcal {H}) \subset \mathbf {Z}.
\]

\begin{enumerate}
\def\labelenumi{\arabic{enumi})}
\setcounter{enumi}{3}
\tightlist
\item
  If \(\mathfrak{g}\) is semi-simple and \(\mathcal{H} = \mathrm{Q}(\mathrm{R}^{\vee})\) , conditions (i) and (ii) of Def. 3 are equivalent, by Th. 3 (ii), to
\end{enumerate}

\begin{enumerate}
\def\labelenumi{(\roman{enumi})}
\setcounter{enumi}{2}
\tightlist
\item
  \(\mathcal{E}\) is stable under \(x^{(n)}\) for all \(\alpha \in \mathrm{R}, x \in \mathcal{G}^{\alpha}, n \in \mathbf{N}\) .
\end{enumerate}

\begin{enumerate}
\def\labelenumi{\arabic{enumi})}
\setcounter{enumi}{4}
\tightlist
\item
  Let B be a basis of R; in conditions (i) and (ii) above, `` \(\alpha \in R\) '' can be replaced by `` \(\alpha \in B \cup (-B)\) '' (loc. cit).
\end{enumerate}

THEOREM 4. Let \(\mathbf{E}\) be a finite dimensional \(\mathfrak{g}\) -module. The following conditions are equivalent:

\begin{enumerate}
\def\labelenumi{(\roman{enumi})}
\item
  E has an admissible lattice;
\item
  every element of \(\mathcal{H}\) defines a diagonalizable endomorphism of \(E\) with integer eigenvalues.
\item
  \(\Longrightarrow\) (ii): this follows from Remark 3.
\item
  \(\Longrightarrow\) (i): we assume that condition (ii) is satisfied and prove (i). By Th. 4 of Chap. I, §6, no. 5, we can assume that the elements of \(\mathfrak{c}\) define homotheties of E, and that E is a simple \(\mathcal{D}\mathfrak{g}\)-module. Let B be a basis of \(\mathrm{R}\), and \(\mathfrak{g} = \mathfrak{n}_{-} \oplus \mathfrak{h} \oplus \mathfrak{n}_{+}\) the corresponding decomposition of \(\mathfrak{g}\). Let \(\lambda\) be the highest weight of the \(\mathcal{D}\mathfrak{g}\)-module E, and let \(e \in E^{\lambda} - \{0\}\) . Put \(\mathcal{E} = \mathcal{U}.e\). It is clear that \(\mathcal{U}.\mathcal{E} \subset \mathcal{E}\) . Since E is simple, \(\mathrm{U}(\mathfrak{g})\) . e = E and hence \(\mathcal{E}\) generates E as a \(\mathbf{Q}\)-vector space. For \(h \in \mathcal{H}\) and \(n \in \mathbf{N}\) , we have \(\binom{h}{n} e = \binom{\lambda(h)}{n} e \in \mathbf{Z} e\) , so
\end{enumerate}

\[
(\mathcal {U} \cap \mathrm{U} (\mathfrak {h})). e = \mathbf {Z}   e.
\]

Since \(\mathrm{U}(\mathfrak{n}_{+}).e=0\) , we have \(\mathcal{E}=(\mathcal{U}\cap\mathrm{U}(\mathfrak{n}_{-}))\) .e by Prop. 3. It now follows from Th. 3 (iii) that \(\mathcal{E}\) is a \(\mathbf{Z}\)-module of finite type.

COROLLARY. If \(\mathfrak{g}\) is semi-simple and \(\mathcal{H}=\mathrm{Q}(\mathrm{R}^{\vee})\) , every finite dimensional \(\mathfrak{g}\)-module has an admissible lattice.

